% Define the hidden glossary
\newacronym{ci}{CI}{Integración Continua}
\newacronym{cd}{CD}{Entrega Continua}
\newacronym{mlops}{MLOps}{Machine Learning Operations}
\newacronym{ml}{ML}{Machine Learning}
\newacronym{ia}{IA}{Inteligencia Artifical}
\newacronym{api}{API}{Interfaz de Programación de Aplicaciones}
\newacronym{cli}{CLI}{Interfaz de Línea de Comandos}
\newacronym{yaml}{YAML}{Yet Another Markup Language}
\newacronym{json}{JSON}{JavaScript Object Notation}
\newacronym{vsc}{VSC}{Visual Studio Code}
\newacronym{hpo}{HPO}{Optimización de Hiperparámetros}
\newacronym{pr}{PR}{Pull Request}
\newacronym{cors}{CORS}{Cross-Origin Resource Sharing}
\newacronym{ui}{UI}{Interfaz de Usuario}
\newacronym{es}{ESM}{Módulos ECMAScript}
\newacronym{hmr}{HMR}{Hot Module Replacement}
\newacronym{dom}{DOM}{Document Object Model}
\newacronym{ide}{IDE}{Integrated Development Environment}
\newacronym{css}{CSS}{Cascading Style Sheets}
\newacronym{html}{HTML}{HyperText Markup Language}
\newacronym{nfs}{NFS}{Sistemas de Archivos en Red}

\newglossaryentry{snapperml}
{
    name=SnapperML,
    description={Framework de \Gls{mlopsg} diseñado para facilitar la reproducibilidad, el seguimiento y la optimización de experimentos de machine learning}
}

\newglossaryentry{mlopsg}{
    name=MLOps,
    description={Conjunto de prácticas que combina el desarrollo y operaciones de software (\Gls{devops}) con el aprendizaje automático, enfocándose en la automatización y gestión eficiente del ciclo de vida de los modelos de machine learning}
}

\newglossaryentry{devops}{
    name=DevOps,
    description={Conjunto de prácticas que combina el desarrollo de software (Development) y las operaciones de TI (Operations) para automatizar y optimizar los procesos de entrega y gestión de aplicaciones, promoviendo la colaboración y la integración continua}
}

\newglossaryentry{gitops}{
    name=GitOps,
    description={Marco operativo basado en prácticas de \Gls{devops}, como la integración continua/entrega continua (\Gls{ci}/\Gls{cd}) y el control de versiones, que automatiza la infraestructura y gestiona la implementación de software}
}

\newglossaryentry{cig}{
    name = CI,
    description = {Práctica de desarrollo en la que el código se integra frecuentemente en un repositorio compartido, seguido de pruebas automáticas, para detectar y corregir errores rápidamente}
}

\newglossaryentry{cdg}{
    name = CD,
    description = {Práctica de automatizar el lanzamiento de cambios de código en producción, asegurando que las actualizaciones sean rápidas, seguras y constantes}
}

\newglossaryentry{deuda}{
    name = {Deuda Técnica},
    text = {deuda técnica},
    description = {Costo a futuro de hacer un desarrollo rápido en lugar de hacerlo de forma correcta, lo que puede generar problemas de mantenimiento y calidad}
}

\newglossaryentry{deudaCapital}{
    name={Deuda Técnica},
    text={Deuda Técnica},
    description={Costo a futuro de hacer un desarrollo rápido en lugar de hacerlo de forma correcta, lo que puede generar problemas de mantenimiento y calidad},
    type={hidden}
}

\newglossaryentry{framework}{
    name = Framework,
    text = framework,
    description = {Conjunto de herramientas, bibliotecas y buenas prácticas que proporcionan una estructura estándar para desarrollar software, facilitando y agilizando el proceso de programación}
}

\newglossaryentry{log}{
    name = Log,
    text = log,
    plural = logs,
    description = {Registro o archivo que almacena información detallada sobre eventos y actividades que ocurren en un sistema o aplicación, como errores, advertencias, o acciones realizadas, para ayudar en el monitoreo, depuración y análisis del rendimiento}
}

\newglossaryentry{apig}{
    name = API,
    plural = APIs,
    description = {Conjunto de reglas y definiciones que permiten a diferentes aplicaciones o sistemas comunicarse entre sí, facilitando el intercambio de datos y la ejecución de funciones}
}

\newglossaryentry{clig}{
    name = CLI,
    description = {Forma de interactuar con un sistema operativo o software mediante la introducción de comandos de texto en lugar de utilizar una interfaz gráfica}
}

\newglossaryentry{hpog}{
    name = HPO,
    description = {Proceso de ajustar los hiperparámetros de un modelo de aprendizaje automático para mejorar su rendimiento. A diferencia de los parámetros aprendidos durante el entrenamiento, los hiperparámetros son configuraciones definidas antes de que el modelo comience a entrenar, como la tasa de aprendizaje o el número de capas en una red neuronal}
}

\newglossaryentry{corsg}{
    name = CORS,
    description = {Mecanismo de seguridad que permite o bloquea el acceso a recursos desde diferentes dominios, controlando cómo los servidores y los navegadores manejan las solicitudes de origen cruzado}
}

\newglossaryentry{ciclomatica}{
    name = {Complejidad Ciclomática},
    text = {complejidad ciclomática},
    description = {Métrica que mide la cantidad de rutas independientes a través del código fuente de un programa. Se utiliza para evaluar la complejidad y el riesgo asociado a la prueba y mantenimiento del software, indicando cuántos casos de prueba mínimos se necesitan para cubrir todas las posibles rutas de ejecución}
}

\newglossaryentry{code smells}{
    name = {Code Smells},
    text = {code smells},
    description = {Indicios en el código que sugieren problemas de diseño o implementación que podrían dificultar su mantenimiento o evolución. No suelen ser errores que impidan que el programa funcione, pero pueden señalar áreas del código que necesitan refactorización para mejorar su calidad}
}

\newglossaryentry{deprecado}{
    name = {Código deprecado},
    description = {Código que aún funciona pero se ha marcado como obsoleto y no se recomienda usarlo, generalmente porque ha sido reemplazado por una mejor alternativa. Puede ser eliminado en futuras versiones, por lo que es importante actualizarlo}
}

\newglossaryentry{artefactos}{
    name = Artefactos,
    text = artefactos,
    description = {Elementos generados y utilizados durante el desarrollo de modelos de aprendizaje automático, como conjuntos de datos procesados, modelos entrenados, métricas de rendimiento, gráficos, y archivos de configuración. Estos artefactos ayudan a documentar y reproducir experimentos}
}

\newglossaryentry{estocasticidad}{
    name = {Estocasticidad},
    text = {estocasticidad},
    description = {Aleatoriedad o incertidumbre inherente en un proceso o sistema. En matemáticas y ciencia de datos, describe fenómenos que no son completamente deterministas, ya que su comportamiento depende de variables aleatorias}
}

\newglossaryentry{frontend}{
    name = Frontend,
    text = frontend,
    description = {Parte del desarrollo web que se ocupa de la interfaz de usuario y la experiencia visual de una aplicación. Incluye todo lo que los usuarios ven e interactúan directamente, como botones, diseño, y contenido, utilizando tecnologías como HTML, CSS y JavaScript}
}

\newglossaryentry{backend}{
    name = Backend,
    text = backend,
    description = {Parte del desarrollo web que gestiona el funcionamiento interno de una aplicación, incluyendo la lógica del negocio, la gestión de bases de datos y la comunicación con el frontend. Se encarga de procesar solicitudes, ejecutar scripts, y almacenar y recuperar datos, utilizando lenguajes como Python, Java, o Node.js}
}

\newglossaryentry{esg}{
    name = {Modulos ES},
    text = {modulos ES},
    description = {Forma de estructurar y organizar el código JavaScript mediante la importación y exportación de funciones, objetos o variables. Facilitan la reutilización y mantenimiento del código, y son soportados de manera nativa por los navegadores modernos}
}

\newglossaryentry{hmrg}{
    name = HMR,
    description = {Funcionalidad en los entornos de desarrollo web que permite actualizar módulos en una aplicación en tiempo real sin necesidad de recargar toda la página}
}

\newglossaryentry{domg}{
    name = {DOM},
    description = {Representación estructurada de un documento HTML o XML en forma de árbol, que permite a los lenguajes de programación manipular y actualizar dinámicamente el contenido, estructura y estilo de la página web.}
}

\newglossaryentry{ideg}{
    name = IDE,
    description = {Entorno de desarrollo que proporciona herramientas integradas como un editor de código, depurador, y compilador para facilitar la creación de software. Ejemplos comunes incluyen Visual Studio Code, PyCharm, y Eclipse.}
}